%% LyX 2.4.4 created this file.  For more info, see https://www.lyx.org/.
%% Do not edit unless you really know what you are doing.
\documentclass[english]{article}
\usepackage[T1]{fontenc}
\usepackage[latin9]{inputenc}
\usepackage{enumitem}
\usepackage{amsmath}
\usepackage{amsthm}
\usepackage{geometry}
\geometry{verbose,tmargin=1in,bmargin=1in,lmargin=1in,rmargin=1in}

\makeatletter
%%%%%%%%%%%%%%%%%%%%%%%%%%%%%% Textclass specific LaTeX commands.
\numberwithin{equation}{section}
\numberwithin{figure}{section}
\newlength{\lyxlabelwidth}      % auxiliary length 

%%%%%%%%%%%%%%%%%%%%%%%%%%%%%% User specified LaTeX commands.
\usepackage{fancyhdr}
\usepackage{lastpage}
\pagestyle{fancy}
\usepackage{enumitem}
\usepackage{ifthen}
%sets math fonts to be more readable
\everymath=\expandafter{\the\everymath\displaystyle}
%\DeclareMathSizes{10}{10}{10}{10}
\usepackage{multicol}
% Added by lyx2lyx
\setlength{\parskip}{\smallskipamount}
\setlength{\parindent}{0pt}

\makeatother

\usepackage{babel}
\begin{document}
\renewcommand{\author}{Dr.\,James Ashe}

\newcommand{\classNum}{331}

\newcommand{\classSec}{A}

\newcommand{\nameType}{}

\newcommand{\examType}{Exam 4 Review}

\renewcommand{\date}{}

%%First page's header%%%%%%%%%%%%%%%%%%%%%%
\fancypagestyle{firststyle} %%header settings for first pagr
{
	\fancyhead[L]{MTH \classNum{}(\classSec{}) \author{}\\
	\textbf{\examType{}} \date{}}
	\fancyhead[C]{\textbf{\nameType}}
	\setlength{\headsep}{14pt}
}
%%%%%%%%%%%%%%%%%%%%%%%%%%%%%%%%%%%%%%%%%%

%%Next page's header%%%%%%%%%%%%%%%%%%%%%%%
\setlist{nolistsep}
\lhead{\classNum{}(\classSec{})} 
\chead{\textbf{\nameType}} 
\cfoot{\thepage{}~of~\pageref{LastPage}} 
\setlength{\headsep}{5pt} 
%%%%%%%%%%%%%%%%%%%%%%%%%%%%%%%%%%%%%%%%%%%

%%Various settings; see the preamble for other settings%%%%%
%%\setenumerate[0]{leftmargin=12pt,itemindent=0pt,labelwidth=0pt}
\renewcommand{\labelenumi}{\textbf{\arabic{enumi}.}}
\renewcommand{\labelenumii}{\textbf{\alph{enumii}.}}
\thispagestyle{firststyle}
%%%%%%%%%%%%%%%%%%%%%%%%%%%%%%%%%%%%%%%%%%

\subsubsection*{13.3 Limits and Continuity}
\begin{itemize}
\item Limits of two-variable continuous functions.

\begin{itemize}
\item \textbf{HW 11, 15; ex. ${\displaystyle \lim_{(x,y)\rightarrow(2,0)}\left(xy^{5}-x\right)}$}
\end{itemize}
\item Limits at the boundary points of two-variable functions. 

\begin{itemize}
\item \textbf{HW 19, 23; ex. ${\displaystyle \lim_{(x,y)\rightarrow(2,0)}\frac{yx-y}{y(x^{2}-1)}}$}
\end{itemize}
\item Nonexistence of limits using the Two-Path Test.

\begin{itemize}
\item \textbf{HW 25, 29; }

\begin{itemize}
\item []\textbf{ex. ${\displaystyle \lim_{(x,y)\rightarrow\left(0,0\right)}\frac{x+2y}{x-2y}}$},
use lines $x=0$ and then $y=0$.
\item []\textbf{ex. ${\displaystyle \lim_{(x,y)\rightarrow\left(0,0\right)}\frac{y^{3}+x^{3}}{xy^{2}}}$},
use lines $x=y$ and then $x=-y$.\vfill
\end{itemize}
\end{itemize}
\end{itemize}

\subsubsection*{13.4 Partial Derivatives}
\begin{itemize}
\item Given a functions $f(x,y)$, find $f_{x}$, $f_{y}$, $f_{xx}$, $f_{yy}$,
$f_{xy}$, and $f_{yx}$. 

\begin{itemize}
\item \textbf{HW 7-15, 17, 23, 27, 29.}\vfill
\end{itemize}
\end{itemize}

\subsubsection*{13.5 The Chain Rule}
\begin{itemize}
\item Use the Chain rule to find the derivatives of functions with one independent
variable. 

\begin{itemize}
\item \textbf{HW 7-10. }(only two-variable functions)\vfill
\end{itemize}
\end{itemize}

\subsubsection*{13.6 Directional Derivatives and the Gradient}
\begin{itemize}
\item Compute the gradient of two-variable functions.

\begin{itemize}
\item \textbf{HW 11, 13.}\vfill
\end{itemize}
\item Find a unit vector that gives the direction of the steepest ascent
and descent of a function.
\item Find a vector that gives a direction of no change of a function.

\begin{itemize}
\item \textbf{HW 21, 25.}\vfill
\end{itemize}
\end{itemize}

\subsubsection*{13.8 Maximum and Minimum Problems}
\begin{itemize}
\item Find the critical points of a function, and use the Second Derivative
Test to determine (when possible) whether each critical point corresponds
to a local maximum, minimum, or saddle point.

\begin{itemize}
\item \textbf{HW 15, 21, 27.}\vfill
\end{itemize}
\end{itemize}

\end{document}
